% This must be in the first 5 lines to tell arXiv to use pdfLaTeX, which is strongly recommended.
\pdfoutput=1
% In particular, the hyperref package requires pdfLaTeX in order to break URLs across lines.

\documentclass[11pt]{article}

% Change "review" to "final" to generate the final (sometimes called camera-ready) version.
% Change to "preprint" to generate a non-anonymous version with page numbers.
\usepackage[review]{acl}

% Standard package includes
\usepackage{times}
\usepackage{latexsym}

% For proper rendering and hyphenation of words containing Latin characters (including in bib files)
\usepackage[T1]{fontenc}
% For Vietnamese characters
% \usepackage[T5]{fontenc}
% See https://www.latex-project.org/help/documentation/encguide.pdf for other character sets

% This assumes your files are encoded as UTF8
\usepackage[utf8]{inputenc}

% This is not strictly necessary, and may be commented out,
% but it will improve the layout of the manuscript,
% and will typically save some space.
\usepackage{microtype}

% This is also not strictly necessary, and may be commented out.
% However, it will improve the aesthetics of text in
% the typewriter font.
\usepackage{inconsolata}

%Including images in your LaTeX document requires adding
%additional package(s)
\usepackage{graphicx}

% If the title and author information does not fit in the area allocated, uncomment the following
%
%\setlength\titlebox{<dim>}
%
% and set <dim> to something 5cm or larger.

\title{SemEval System Papers}
\author{Santiago Lares Harbin \\
  Universidad de San Andres \\
  \texttt{santiago.lares@udesa.edu.ar} \\\And
  Team: Shouth NLP \\
  \texttt{team@shouthnlp.com} \\}
\date{Santiago Lares Harbin \\ Universidad de San Andres \\ \texttt{santiago.lares@udesa.edu.ar} \\ Team: Shouth NLP \\ \texttt{team@shouthnlp.com}}


\author{Santiago Lares Harbin \\
  Universidad de San Andres \\
  \texttt{santiago.lares@udesa.edu.ar} \\\And
  Team: Shouth NLP \\
  \texttt{team@shouthnlp.com} \\}

%\author{
%  \textbf{First Author\textsuperscript{1}},
%  \textbf{Second Author\textsuperscript{1,2}},
%  \textbf{Third T. Author\textsuperscript{1}},
%  \textbf{Fourth Author\textsuperscript{1}},
%\\
%  \textbf{Fifth Author\textsuperscript{1,2}},
%  \textbf{Sixth Author\textsuperscript{1}},
%  \textbf{Seventh Author\textsuperscript{1}},
%  \textbf{Eighth Author \textsuperscript{1,2,3,4}},
%\\
%  \textbf{Ninth Author\textsuperscript{1}},
%  \textbf{Tenth Author\textsuperscript{1}},
%  \textbf{Eleventh E. Author\textsuperscript{1,2,3,4,5}},
%  \textbf{Twelfth Author\textsuperscript{1}},
%\\
%  \textbf{Thirteenth Author\textsuperscript{3}},
%  \textbf{Fourteenth F. Author\textsuperscript{2,4}},
%  \textbf{Fifteenth Author\textsuperscript{1}},
%  \textbf{Sixteenth Author\textsuperscript{1}},
%\\
%  \textbf{Seventeenth S. Author\textsuperscript{4,5}},
%  \textbf{Eighteenth Author\textsuperscript{3,4}},
%  \textbf{Nineteenth N. Author\textsuperscript{2,5}},
%  \textbf{Twentieth Author\textsuperscript{1}}
%\\
%\\
%  \textsuperscript{1}Affiliation 1,
%  \textsuperscript{2}Affiliation 2,
%  \textsuperscript{3}Affiliation 3,
%  \textsuperscript{4}Affiliation 4,
%  \textsuperscript{5}Affiliation 5
%\\
%  \small{
%    \textbf{Correspondence:} \href{mailto:email@domain}{email@domain}
%  }
%}

\begin{document}
\maketitle
\begin{abstract}
Abstract
\end{abstract}

\section{Introduction}

Introduction

\cite{Gusfield:97}

\section{Background}

Background

\section{System Overview}

Our system addresses the challenge of retrieving relevant fact checks from a large corpus of 200,000 documents by implementing a fine-tuned contrastive learning approach based on the multilingual E5 model architecture. The system consists of three main components: (1) a text embedding model, (2) a contrastive learning framework, and (3) a retrieval pipeline.

\subsection{Model Architecture}

The core of our system is built upon the `intfloat/multilingual-e5-large-instruct` model, which we enhance with a custom architecture designed specifically for fact-checking retrieval. The model consists of:

\begin{itemize}
    \item A transformer encoder (base E5 model)
    \item A mean pooling layer for sentence representation
    \item A dense projection layer with GELU activation
    \item L2 normalization for cosine similarity optimization
\end{itemize}

The model processes both posts and fact checks using language-specific prefixes ('query:' and 'passage:' respectively) to optimize the semantic matching between them.

\subsection{Contrastive Learning Framework}

We implement a contrastive learning approach using the Multiple Negatives Ranking Loss (MNRL), which can be formulated as:

\begin{equation}
L = -\log \frac{\exp(sim(x_i, x_i^+)/\tau)}{\sum_{j=1}^{N} \exp(sim(x_i, x_j)/\tau)}
\end{equation}

where:
\begin{itemize}
    \item $x_i$ represents the anchor text embedding (post)
    \item $x_i^+$ represents the positive pair embedding (matching fact-check)
    \item $x_j$ represents all other samples in the batch (negative pairs)
    \item $\tau$ is a temperature parameter that controls the sharpness of the distribution
    \item $sim(\cdot,\cdot)$ is the cosine similarity function
\end{itemize}

The framework employs in-batch negatives for computational efficiency while maintaining effective contrastive learning. The temperature parameter $\tau$ is dynamically adjusted during training to optimize the learning process.

\subsection{Retrieval Pipeline}

The retrieval process follows these steps:

\begin{enumerate}
    \item \textbf{Language Detection}: The system first identifies the language of the input post.
    \item \textbf{Corpus Filtering}: The fact-check corpus is filtered to prioritize documents in the same language as the input post.
    \item \textbf{Embedding Generation}: The model generates embeddings for both the input post and the filtered fact checks.
    \item \textbf{Similarity Ranking}: The system computes Pearson cosine similarity between the post and fact check embeddings.
    \item \textbf{Top-K Selection}: The 10 most similar fact checks are retrieved and ranked.
\end{enumerate}

This pipeline is optimized for both monolingual and cross-lingual retrieval, with special consideration for language-specific nuances in the embedding space.

\section{Experimental Setup}

\subsection{Dataset}

The experimental data consists of three main components:
\begin{itemize}
    \item \textbf{Posts}: Social media posts in multiple languages requiring fact-checking
    \item \textbf{Fact Checks}: A corpus of 153,743 fact-checks across 8 languages (Arabic, German, English, French, Malaysian, Portuguese, Spanish, and Thai)
    \item \textbf{Training Pairs}: Annotated pairs of posts and their corresponding relevant fact-checks
\end{itemize}

We split the training data into training (80\%) and validation (20\%) sets while maintaining the language distribution. For each post, we concatenated the OCR-extracted text and manual transcription to create a comprehensive text representation. Similarly, for fact checks, we combined the claim and title fields to create rich semantic representations.

\subsection{Implementation Details}

The system was implemented using the following configuration:

\begin{itemize}
    \item \textbf{Base Model}: multilingual-e5-large-instruct
    \item \textbf{Framework}: SentenceTransformers (v3.4.1)
    \item \textbf{Deep Learning Backend}: PyTorch (v2.5.1)
    \item \textbf{Python Version}: 3.12.7
    \item \textbf{Computing Environment}: macOS 15.1 (24B83)
\end{itemize}

\subsection{Training Configuration}

The model was trained with the following hyperparameters:

\begin{table}[h]
\centering
\small
\begin{tabular}{ll}
\hline
\textbf{Parameter} & \textbf{Value} \\
\hline
Maximum Sequence Length & 512 tokens \\
Pooling Mode & Mean \\
Embedding Dimension & 384 \\
Batch Size & 4 \\
Evaluation Batch Size & 16 \\
Gradient Accumulation Steps & 2 \\
Learning Rate & 2e-5 \\
Temperature & 0.05 \\
Mixed Precision Training & Enabled \\
\hline
\end{tabular}
\caption{Model training hyperparameters}
\label{tab:hyperparams}
\end{table}

\subsection{Evaluation Metrics}

The primary evaluation metric for our system was Success@10 (S@10), which measures the proportion of posts for which the correct fact check appears in the top 10 retrieved results. We evaluated the model's performance in both monolingual and cross-lingual settings:

\begin{itemize}
    \item \textbf{Monolingual}: Retrieval within the same language
    \item \textbf{Cross-lingual}: Retrieval across different languages
\end{itemize}

During training, we also monitored:
\begin{itemize}
    \item Cosine similarity between positive pairs
    \item Cosine similarity between negative pairs
    \item Training loss convergence
    \item Validation performance at regular intervals
\end{itemize}

\section{Results}

We evaluate our system's performance on both the validation and test sets, analyzing the impact of different model configurations and language-specific performance. Our system achieved an average Success@10 score of 0.867 on the official test set, ranking 20th out of 28 published participants.

\subsection{Model Configuration Analysis}

Table~\ref{tab:base-models} shows the performance comparison of different base models and input prefix configurations on the validation set. The results demonstrate that the multilingual-e5-large-instruct model with appropriate prefixes consistently outperforms other configurations.

\begin{table*}[t]
  \centering
  \small
  \begin{tabular}{p{5cm}p{3cm}p{3cm}c}
    \hline
    \textbf{Model} & \textbf{Posts Prefix} & \textbf{Facts Prefix} & \textbf{Success@10} \\
    \hline
    multilingual-e5-large-instruct & passage: & query: & \textbf{0.834} \\
    multilingual-e5-large-instruct & -- & -- & 0.821 \\
    multilingual-e5-small & -- & -- & 0.795 \\
    e5-large-v2 & -- & -- & 0.758 \\
    modernbert-embed-base & search\_query: & search\_document: & 0.779 \\
    paraphrase-multilingual-mpnet-base-v2 & -- & -- & 0.736 \\
    \hline
  \end{tabular}
  \caption{Performance comparison of base models on the validation set. Models are ordered by Success@10 score. Best result is shown in \textbf{bold}.}
  \label{tab:base-models}
\end{table*}

\subsection{Language-Specific Performance}

Table~\ref{tab:lang-performance} presents the detailed performance breakdown by language on the validation set, showing significant variation across different languages.

\begin{table*}[t]
  \centering
  \small
  \begin{tabular}{lccc}
    \hline
    \textbf{Language} & \textbf{Success@10} & \textbf{Sample Count} & \textbf{Relative Performance} \\
    \hline
    Thai (tha) & \textbf{0.930} & 372 & +10.9\% \\
    French (fra) & 0.895 & 1,277 & +7.4\% \\
    Spanish (spa) & 0.857 & 4,502 & +3.6\% \\
    Arabic (ara) & 0.830 & 541 & +0.9\% \\
    Malaysian (msa) & 0.817 & 849 & -0.4\% \\
    German (deu) & 0.777 & 533 & -4.4\% \\
    Portuguese (por) & 0.781 & 2,057 & -4.0\% \\
    English (eng) & 0.766 & 3,481 & -5.5\% \\
    \hline
    \textbf{Average} & 0.821 & 13,612 & -- \\
    \hline
  \end{tabular}
  \caption{Language-specific performance on the validation set. Relative Performance shows the percentage difference from the average Success@10 score. Best performance is shown in \textbf{bold}.}
  \label{tab:lang-performance}
\end{table*}

\subsection{Impact of Input Prefixes}

Our experiments revealed that the choice of input prefixes significantly impacts model performance. Using the recommended prefixes ('passage:' and 'query:') improved the average Success@10 score by 0.013 points (from 0.821 to 0.834) compared to using no prefixes. This improvement was consistent across all languages, with the largest gains observed in:

\begin{itemize}
    \item English: +0.022 (from 0.766 to 0.788)
    \item Spanish: +0.012 (from 0.857 to 0.869)
    \item Portuguese: +0.018 (from 0.781 to 0.799)
\end{itemize}

These results suggest that appropriate input formatting plays a crucial role in optimizing the model's cross-lingual understanding and retrieval capabilities.

\section{Conclusion}

Conclusion


\section{Acknowledgments}

Acknowledgments

% Bibliography entries for the entire Anthology, followed by custom entries
%\bibliography{anthology,custom}
% Custom bibliography entries only
\bibliography{custom}

\appendix

\section{Detailed Results by Language}
\label{sec:appendix-results}

\begin{table*}[t]
  \centering
  \small
  \begin{tabular}{lcc}
    \hline
    \textbf{Language} & \textbf{Success@10} & \textbf{Test Set Size} \\
    \hline
    Malaysian (msa) & \textbf{0.978} & -- \\
    Thai (tha) & 0.940 & -- \\
    Arabic (ara) & 0.912 & -- \\
    French (fra) & 0.894 & -- \\
    Spanish (spa) & 0.866 & -- \\
    German (deu) & 0.858 & -- \\
    Turkish (tur) & 0.828 & -- \\
    Polish (pol) & 0.806 & -- \\
    English (eng) & 0.796 & -- \\
    Portuguese (por) & 0.796 & -- \\
    \hline
    Average & 0.867 & -- \\
    \hline
  \end{tabular}
  \caption{Official test set results of the fine-tuned multilingual-e5-large-instruct model by language. Languages are ordered by Success@10 score. Best performance is shown in \textbf{bold}. The model shows strong cross-lingual generalization, with particularly strong performance on Malaysian and Thai.}
  \label{tab:test-results-by-language}
\end{table*}

This appendix presents the detailed performance breakdown of our fine-tuned model across different languages in the official test set. As shown in Table~\ref{tab:test-results-by-language}, the model achieves consistent performance across diverse language families, from Indo-European (English, French, German) to Austronesian (Malaysian) and Tai-Kadai (Thai) languages.

\end{document}
